\subsection{小批次改变竞争区间}
\label{sec:7-smallbatch}

T5 在四个 E1 数据集、batch 8/32 下比较调度 JHQ、CAGRA fp32/int8 和 IVF-RaBitQ。
五轮同步主机到主机调用分别使用 1,000/992 条查询。1,808 个点包含 JHQ 单 block
对照；本实验不含 IVF-PQ 和暴力搜索。
\Cref{fig:7-smallbatch,tab:7-smallbatch}比较所得到的工作区间。

\suppressfloats[t]
\begin{figure}[htbp]
\centering
\includegraphics[width=\linewidth]{t4-t5/smallbatch.pdf}
\caption[小批次竞争区间]{T5 包含达标实测端点的插值估计下，JHQ／最快对手 QPS。
大于一表示 JHQ 更快。严格实测阈值见\cref{tab:7-smallbatch}。}
\label{fig:7-smallbatch}
\end{figure}

\begin{table}[htbp]
\centering
\caption[小批次同召回率吞吐比]{T5 QPS 比值，依次为“包含达标端点的插值／严格实测阈值”。
每种估计独立选择最快对手；更高 recall 的达标端点保留，无外推。}
\label{tab:7-smallbatch}
\footnotesize\setlength{\tabcolsep}{4pt}
\begin{tabular}{@{}lrrrr@{}}
\toprule
数据集 & $b$ & $R=0.90$ & $R=0.95$ & $R=0.97$ \\
\midrule
Vogue & 8 & 1.155 / 1.091 & 1.073 / 0.980 & 0.893 / 0.854 \\
Vogue & 32 & 1.069 / 0.997 & 0.956 / 0.895 & 0.865 / 0.839 \\
arXiv & 8 & 1.479 / 1.454 & 0.859 / 0.851 & 0.697 / 0.698 \\
arXiv & 32 & 1.166 / 1.148 & 0.667 / 0.661 & 0.622 / 0.582 \\
OpenAI-1536 & 8 & 1.316 / 1.131 & 0.851 / 0.825 & 0.641 / 0.620 \\
OpenAI-1536 & 32 & 1.071 / 0.915 & 0.770 / 0.752 & 0.937 / 0.914 \\
OpenAI-3072 & 8 & 1.032 / 0.953 & 0.798 / 0.798 & 0.721 / 0.702 \\
OpenAI-3072 & 32 & 0.866 / 0.734 & 0.608 / 0.771 & 1.008 / 1.351 \\
\bottomrule
\end{tabular}

\end{table}

recall 0.97、batch 8 下，两种估计均显示 CAGRA fp32 在四个数据集领先，
JHQ 插值比为 0.641--0.893。batch 32 的插值比为 0.622--1.008；
OpenAI-3072 在插值下接近平手，但实测阈值有利于 JHQ。
recall 0.95 时，Vogue 的 batch-8 领先同样只出现在插值估计。
因此，边界排名受网格分辨率影响。

较低 recall 更有利于 JHQ，但必须保留最快前沿点已经超过目标的竞争者。
0.90 处，修正后的插值比为 batch 8 的 1.032--1.479 和 batch 32 的
0.866--1.166；严格实测阈值分别给出三个和一个 JHQ 领先。
排除这些达标端点会夸大低 recall 优势。
调度仍然有用：0.95 处相对单 block 的插值提速为 1.7--3.3$\times$。
内部执行提速因此未必转化为跨系统领先。

按数据集、方法与 batch 汇总的基线空调用中位数为 0.012--0.072\,ms，
该开销仍包含在搜索计时内。
CAGRA 使用 Python 绑定，JHQ／RaBitQ 使用 C++，预热策略也有差别。
证据支持有限网格、单客户端的工作区间，协议见\cref{sec:a-t4-t5}。
E1 在 batch 1,000 的高 recall 领先不能推广到小批次。
